\documentclass[11pt,letterpaper]{article}
\usepackage[technical-reference]{cornell-notes}
\usepackage{needspace}

\title{Clause 13: Templates - Cornell Notes}
\author{}
\date{}
\setDocTitle{Clause 13: Templates}
\setDocSubtitle{C++ 2024 Cornell Notes}
\setDocOwner{C++ 2024 Cornell Notes}
\setDocAuthor{}
\setDocDate{}
\setCornellCollection{C++ 2024 Cornell Notes}
\setCornellUnitType{Clause}
\setCornellUnitNumber{13}
\setCornellUnitTitle{Templates}
\setCornellCompanionLabel{C++ Technical Reference}
\hypersetup{pdftitle={Clause 13: Templates - Cornell Notes},pdfsubject={C++ technical study notes},pdfauthor={}}

\begin{document}
\maketitle
\begin{CornellReferenceOrientationBox}
\textbf{Purpose.} Defines template parameters and arguments, constraints, declaration forms, specialization, dependent names, instantiation, deduction, partial ordering, and overload resolution for templates.
\par\textbf{Role.} Provides compile-time abstraction and controlled generation of declarations and definitions.
\par\textbf{Principal dependencies.} Uses declarations, lookup, expressions, classes, and overloading; concepts and library facilities build on it.
\par\textbf{Primary audience.} programmers, library authors, and compiler implementers.
\par\textbf{Major subject groups.} Preamble; Template parameters; Names of template specializations; Template arguments; Template constraints; Type equivalence; Template declarations; Name resolution; Template instantiation and specialization; Function template specializations.
\end{CornellReferenceOrientationBox}
\section{Learning Objectives}
\begin{itemize}[label=\textcolor{CornellReferenceTeal}{$\blacktriangleright$}]
\item Define the governing ideas of Preamble and state the consequences for a program or implementation.
\item Identify the governing ideas of Template parameters and state the consequences for a program or implementation.
\item Distinguish the governing ideas of Names of template specializations and state the consequences for a program or implementation.
\item Explain the governing ideas of Template arguments and state the consequences for a program or implementation.
\item Trace the governing ideas of Template constraints and state the consequences for a program or implementation.
\item Apply the governing ideas of Type equivalence and state the consequences for a program or implementation.
\item Analyze the governing ideas of Template declarations and state the consequences for a program or implementation.
\item Evaluate the governing ideas of Name resolution and state the consequences for a program or implementation.
\item Diagnose the governing ideas of Template instantiation and specialization and state the consequences for a program or implementation.
\item Compare the governing ideas of Function template specializations and state the consequences for a program or implementation.
\item Verify the governing ideas of General and state the consequences for a program or implementation.
\item Relate the governing ideas of Template type arguments and state the consequences for a program or implementation.
\item Classify the governing ideas of Template non-type arguments and state the consequences for a program or implementation.
\item Justify the governing ideas of Template template arguments and state the consequences for a program or implementation.
\item Audit the governing ideas of General and state the consequences for a program or implementation.
\item Summarize the governing ideas of Constraints and state the consequences for a program or implementation.
\end{itemize}
\section{Normative Structure Map}
\small\begin{longtable}{@{}>{\RaggedRight\arraybackslash}p{0.13\textwidth}>{\RaggedRight\arraybackslash}p{0.27\textwidth}>{\RaggedRight\arraybackslash}p{0.19\textwidth}>{\RaggedRight\arraybackslash}p{0.31\textwidth}@{}}
\toprule\textbf{Subclause} & \textbf{Subject} & \textbf{Rule category} & \textbf{Key dependencies / entities}\\\midrule\endfirsthead
\toprule\textbf{Subclause} & \textbf{Subject} & \textbf{Rule category} & \textbf{Key dependencies / entities}\\\midrule\endhead
\bottomrule\endfoot
13.1 & Preamble & core-language semantics & Uses declarations, lookup, expressions, classes, and overloading; concepts and library facilities build on it.\\
13.2 & Template parameters & core-language semantics & Uses declarations, lookup, expressions, classes, and overloading; concepts and library facilities build on it.\\
13.3 & Names of template specializations & core-language semantics & Uses declarations, lookup, expressions, classes, and overloading; concepts and library facilities build on it.\\
13.4 & Template arguments & core-language semantics & General, Template type arguments, Template non-type arguments, Template template arguments\\
13.5 & Template constraints & requirements & General, Constraints, Constrained declarations, Constraint normalization, and related subclauses\\
13.6 & Type equivalence & core-language semantics & Uses declarations, lookup, expressions, classes, and overloading; concepts and library facilities build on it.\\
13.7 & Template declarations & core-language semantics & General, Class templates, Member templates, Variadic templates, and related subclauses\\
13.8 & Name resolution & core-language semantics & General, Locally declared names, Dependent names, Dependent name resolution\\
13.9 & Template instantiation and specialization & core-language semantics & General, Implicit instantiation, Explicit instantiation, Explicit specialization\\
13.10 & Function template specializations & operation semantics & General, Explicit template argument specification, Template argument deduction, Overload resolution\\
\end{longtable}\normalsize
\begin{CornellReferenceNormativeBox}A heading maps the clause; it does not replace the detailed conditions. For each operation, read requirements, effects, result, exceptions, complexity, remarks, and notes with their stated normative force.\end{CornellReferenceNormativeBox}
\subsection*{Rule Classification Guide}
\small\begin{longtable}{@{}>{\RaggedRight\arraybackslash}p{0.25\textwidth}>{\RaggedRight\arraybackslash}p{0.67\textwidth}@{}}\toprule\textbf{Label or category} & \textbf{How to read it}\\\midrule\endfirsthead\toprule\textbf{Label or category} & \textbf{How to read it (continued)}\\\midrule\endhead\bottomrule\endfoot
Well-formed / ill-formed & Formation is governed by syntax and semantic rules; an ill-formed program does not become well-formed because one implementation accepts it.\\
Diagnosable rule & A violation requires at least one diagnostic unless the wording places the case outside the diagnosable set.\\
No diagnostic required & The program can be ill-formed while the implementation has no diagnostic obligation.\\
Undefined behavior & No execution requirements are imposed; translation success does not establish safety or portability.\\
Unspecified behavior & One of the permitted outcomes may be chosen without documentation.\\
Implementation-defined behavior & The implementation chooses and documents the behavior.\\
Conditionally supported & The implementation may omit support but documents unsupported constructs.\\
\end{longtable}\normalsize
\section{Cornell Cue-and-Notes}
\Needspace{8\baselineskip}
\subsection*{13.1 Preamble}
\begin{longtable}{@{}>{\RaggedRight\arraybackslash\columncolor{CornellReferenceCueGray}}p{0.28\textwidth}>{\RaggedRight\arraybackslash}p{0.64\textwidth}@{}}
\rowcolor{CornellReferenceNavy}\color{white}\textbf{Cue / recall prompt} & \color{white}\textbf{Notes}\\\endfirsthead
\rowcolor{CornellReferenceNavy}\color{white}\textbf{Cue / recall prompt} & \color{white}\textbf{Notes (continued)}\\\endhead
\arrayrulecolor{CornellReferenceLineGray}
\textbf{What rule applies to Preamble?}\\[-1mm]{\scriptsize\CornellClauseRef{13.1}} & Frames the terminology, applicability, and relationships needed to interpret Preamble; detailed rules remain in the following subclauses.\\\hline
\end{longtable}
\begin{CornellReferenceSummaryBox}Preamble supplies the core-language semantics portion of this clause. The key review move is to connect its local wording to the clause-wide model, then classify each condition by normative force before relying on it.\end{CornellReferenceSummaryBox}
\Needspace{8\baselineskip}
\subsection*{13.2 Template parameters}
\begin{longtable}{@{}>{\RaggedRight\arraybackslash\columncolor{CornellReferenceCueGray}}p{0.28\textwidth}>{\RaggedRight\arraybackslash}p{0.64\textwidth}@{}}
\rowcolor{CornellReferenceNavy}\color{white}\textbf{Cue / recall prompt} & \color{white}\textbf{Notes}\\\endfirsthead
\rowcolor{CornellReferenceNavy}\color{white}\textbf{Cue / recall prompt} & \color{white}\textbf{Notes (continued)}\\\endhead
\arrayrulecolor{CornellReferenceLineGray}
\textbf{What rule applies to Template parameters?}\\[-1mm]{\scriptsize\CornellClauseRef{13.2}} & Defines the template context for Template parameters; track dependence, lookup time, substitution, constraints, instantiation, and specialization ordering separately.\\\hline
\end{longtable}
\begin{CornellReferenceSummaryBox}Template parameters supplies the core-language semantics portion of this clause. The key review move is to connect its local wording to the clause-wide model, then classify each condition by normative force before relying on it.\end{CornellReferenceSummaryBox}
\Needspace{8\baselineskip}
\subsection*{13.3 Names of template specializations}
\begin{longtable}{@{}>{\RaggedRight\arraybackslash\columncolor{CornellReferenceCueGray}}p{0.28\textwidth}>{\RaggedRight\arraybackslash}p{0.64\textwidth}@{}}
\rowcolor{CornellReferenceNavy}\color{white}\textbf{Cue / recall prompt} & \color{white}\textbf{Notes}\\\endfirsthead
\rowcolor{CornellReferenceNavy}\color{white}\textbf{Cue / recall prompt} & \color{white}\textbf{Notes (continued)}\\\endhead
\arrayrulecolor{CornellReferenceLineGray}
\textbf{What rule applies to Names of template specializations?}\\[-1mm]{\scriptsize\CornellClauseRef{13.3}} & Defines the template context for Names of template specializations; track dependence, lookup time, substitution, constraints, instantiation, and specialization ordering separately.\\\hline
\end{longtable}
\begin{CornellReferenceSummaryBox}Names of template specializations supplies the core-language semantics portion of this clause. The key review move is to connect its local wording to the clause-wide model, then classify each condition by normative force before relying on it.\end{CornellReferenceSummaryBox}
\Needspace{8\baselineskip}
\subsection*{13.4 Template arguments}
\begin{longtable}{@{}>{\RaggedRight\arraybackslash\columncolor{CornellReferenceCueGray}}p{0.28\textwidth}>{\RaggedRight\arraybackslash}p{0.64\textwidth}@{}}
\rowcolor{CornellReferenceNavy}\color{white}\textbf{Cue / recall prompt} & \color{white}\textbf{Notes}\\\endfirsthead
\rowcolor{CornellReferenceNavy}\color{white}\textbf{Cue / recall prompt} & \color{white}\textbf{Notes (continued)}\\\endhead
\arrayrulecolor{CornellReferenceLineGray}
\textbf{What rule applies to General?}\\[-1mm]{\scriptsize\CornellClauseRef{13.4.1}} & Frames the terminology, applicability, and relationships needed to interpret General; detailed rules remain in the following subclauses.\\\hline
\textbf{What rule applies to Template type arguments?}\\[-1mm]{\scriptsize\CornellClauseRef{13.4.2}} & Defines the template context for Template type arguments; track dependence, lookup time, substitution, constraints, instantiation, and specialization ordering separately.\\\hline
\textbf{What rule applies to Template non-type arguments?}\\[-1mm]{\scriptsize\CornellClauseRef{13.4.3}} & Defines the template context for Template non-type arguments; track dependence, lookup time, substitution, constraints, instantiation, and specialization ordering separately.\\\hline
\textbf{What rule applies to Template template arguments?}\\[-1mm]{\scriptsize\CornellClauseRef{13.4.4}} & Defines the template context for Template template arguments; track dependence, lookup time, substitution, constraints, instantiation, and specialization ordering separately.\\\hline
\end{longtable}
\begin{CornellReferenceSummaryBox}Template arguments supplies the core-language semantics portion of this clause. The key review move is to connect its local wording to the clause-wide model, then classify each condition by normative force before relying on it.\end{CornellReferenceSummaryBox}
\Needspace{8\baselineskip}
\subsection*{13.5 Template constraints}
\begin{longtable}{@{}>{\RaggedRight\arraybackslash\columncolor{CornellReferenceCueGray}}p{0.28\textwidth}>{\RaggedRight\arraybackslash}p{0.64\textwidth}@{}}
\rowcolor{CornellReferenceNavy}\color{white}\textbf{Cue / recall prompt} & \color{white}\textbf{Notes}\\\endfirsthead
\rowcolor{CornellReferenceNavy}\color{white}\textbf{Cue / recall prompt} & \color{white}\textbf{Notes (continued)}\\\endhead
\arrayrulecolor{CornellReferenceLineGray}
\textbf{What rule applies to General?}\\[-1mm]{\scriptsize\CornellClauseRef{13.5.1}} & Frames the terminology, applicability, and relationships needed to interpret General; detailed rules remain in the following subclauses.\\\hline
\textbf{What rule applies to Constraints?}\\[-1mm]{\scriptsize\CornellClauseRef{13.5.2}} & States the admissibility and semantic obligations for Constraints. Separate compile-time participation rules from caller preconditions and runtime effects.\\\hline
\textbf{What rule applies to Constrained declarations?}\\[-1mm]{\scriptsize\CornellClauseRef{13.5.3}} & Defines the core-language rules for Constrained declarations, including formation, semantic effect, interactions with type and lookup, and any ill-formed or implementation-dependent cases.\\\hline
\textbf{What rule applies to Constraint normalization?}\\[-1mm]{\scriptsize\CornellClauseRef{13.5.4}} & States the admissibility and semantic obligations for Constraint normalization. Separate compile-time participation rules from caller preconditions and runtime effects.\\\hline
\textbf{What rule applies to Partial ordering by constraints?}\\[-1mm]{\scriptsize\CornellClauseRef{13.5.5}} & States the admissibility and semantic obligations for Partial ordering by constraints. Separate compile-time participation rules from caller preconditions and runtime effects.\\\hline
\end{longtable}
\begin{CornellReferenceSummaryBox}Template constraints supplies the requirements portion of this clause. The key review move is to connect its local wording to the clause-wide model, then classify each condition by normative force before relying on it.\end{CornellReferenceSummaryBox}
\Needspace{5\baselineskip}\paragraph{Detailed coverage ledger.}
\begin{description}[style=nextline,leftmargin=2.8cm,font=\normalfont\bfseries\color{CornellReferenceTeal}]
\item[13.5.2.1] \textbf{General.} Frames the terminology, applicability, and relationships needed to interpret General; detailed rules remain in the following subclauses.
\item[13.5.2.2] \textbf{Logical operations.} Defines the observable result of Logical operations together with applicable iterator-domain, ordering, complexity, invalidation, and exception conditions.
\item[13.5.2.3] \textbf{Atomic constraints.} States the admissibility and semantic obligations for Atomic constraints. Separate compile-time participation rules from caller preconditions and runtime effects.
\end{description}
\Needspace{8\baselineskip}
\subsection*{13.6 Type equivalence}
\begin{longtable}{@{}>{\RaggedRight\arraybackslash\columncolor{CornellReferenceCueGray}}p{0.28\textwidth}>{\RaggedRight\arraybackslash}p{0.64\textwidth}@{}}
\rowcolor{CornellReferenceNavy}\color{white}\textbf{Cue / recall prompt} & \color{white}\textbf{Notes}\\\endfirsthead
\rowcolor{CornellReferenceNavy}\color{white}\textbf{Cue / recall prompt} & \color{white}\textbf{Notes (continued)}\\\endhead
\arrayrulecolor{CornellReferenceLineGray}
\textbf{What rule applies to Type equivalence?}\\[-1mm]{\scriptsize\CornellClauseRef{13.6}} & Defines the core-language rules for Type equivalence, including formation, semantic effect, interactions with type and lookup, and any ill-formed or implementation-dependent cases.\\\hline
\end{longtable}
\begin{CornellReferenceSummaryBox}Type equivalence supplies the core-language semantics portion of this clause. The key review move is to connect its local wording to the clause-wide model, then classify each condition by normative force before relying on it.\end{CornellReferenceSummaryBox}
\Needspace{8\baselineskip}
\subsection*{13.7 Template declarations}
\begin{longtable}{@{}>{\RaggedRight\arraybackslash\columncolor{CornellReferenceCueGray}}p{0.28\textwidth}>{\RaggedRight\arraybackslash}p{0.64\textwidth}@{}}
\rowcolor{CornellReferenceNavy}\color{white}\textbf{Cue / recall prompt} & \color{white}\textbf{Notes}\\\endfirsthead
\rowcolor{CornellReferenceNavy}\color{white}\textbf{Cue / recall prompt} & \color{white}\textbf{Notes (continued)}\\\endhead
\arrayrulecolor{CornellReferenceLineGray}
\textbf{What rule applies to General?}\\[-1mm]{\scriptsize\CornellClauseRef{13.7.1}} & Frames the terminology, applicability, and relationships needed to interpret General; detailed rules remain in the following subclauses.\\\hline
\textbf{What rule applies to Class templates?}\\[-1mm]{\scriptsize\CornellClauseRef{13.7.2}} & Defines the template context for Class templates; track dependence, lookup time, substitution, constraints, instantiation, and specialization ordering separately.\\\hline
\textbf{What rule applies to Member templates?}\\[-1mm]{\scriptsize\CornellClauseRef{13.7.3}} & Defines the template context for Member templates; track dependence, lookup time, substitution, constraints, instantiation, and specialization ordering separately.\\\hline
\textbf{What rule applies to Variadic templates?}\\[-1mm]{\scriptsize\CornellClauseRef{13.7.4}} & Defines the template context for Variadic templates; track dependence, lookup time, substitution, constraints, instantiation, and specialization ordering separately.\\\hline
\textbf{What rule applies to Friends?}\\[-1mm]{\scriptsize\CornellClauseRef{13.7.5}} & Defines the core-language rules for Friends, including formation, semantic effect, interactions with type and lookup, and any ill-formed or implementation-dependent cases.\\\hline
\textbf{What rule applies to Partial specialization?}\\[-1mm]{\scriptsize\CornellClauseRef{13.7.6}} & Defines the template context for Partial specialization; track dependence, lookup time, substitution, constraints, instantiation, and specialization ordering separately.\\\hline
\textbf{What rule applies to Function templates?}\\[-1mm]{\scriptsize\CornellClauseRef{13.7.7}} & Defines the template context for Function templates; track dependence, lookup time, substitution, constraints, instantiation, and specialization ordering separately.\\\hline
\textbf{What rule applies to Alias templates?}\\[-1mm]{\scriptsize\CornellClauseRef{13.7.8}} & Defines the template context for Alias templates; track dependence, lookup time, substitution, constraints, instantiation, and specialization ordering separately.\\\hline
\textbf{What rule applies to Concept definitions?}\\[-1mm]{\scriptsize\CornellClauseRef{13.7.9}} & Defines or applies the generic requirement Concept definitions; syntactic satisfaction must be paired with every stated semantic and equality-preservation expectation.\\\hline
\end{longtable}
\begin{CornellReferenceSummaryBox}Template declarations supplies the core-language semantics portion of this clause. The key review move is to connect its local wording to the clause-wide model, then classify each condition by normative force before relying on it.\end{CornellReferenceSummaryBox}
\Needspace{5\baselineskip}\paragraph{Detailed coverage ledger.}
\begin{description}[style=nextline,leftmargin=2.8cm,font=\normalfont\bfseries\color{CornellReferenceTeal}]
\item[13.7.2.1] \textbf{General.} Frames the terminology, applicability, and relationships needed to interpret General; detailed rules remain in the following subclauses.
\item[13.7.2.2] \textbf{Member functions of class templates.} Defines the template context for Member functions of class templates; track dependence, lookup time, substitution, constraints, instantiation, and specialization ordering separately.
\item[13.7.2.3] \textbf{Deduction guides.} Defines the template context for Deduction guides; track dependence, lookup time, substitution, constraints, instantiation, and specialization ordering separately.
\item[13.7.2.4] \textbf{Member classes of class templates.} Defines the template context for Member classes of class templates; track dependence, lookup time, substitution, constraints, instantiation, and specialization ordering separately.
\item[13.7.2.5] \textbf{Static data members of class templates.} Defines the template context for Static data members of class templates; track dependence, lookup time, substitution, constraints, instantiation, and specialization ordering separately.
\item[13.7.2.6] \textbf{Enumeration members of class templates.} Defines the template context for Enumeration members of class templates; track dependence, lookup time, substitution, constraints, instantiation, and specialization ordering separately.
\item[13.7.6.1] \textbf{General.} Frames the terminology, applicability, and relationships needed to interpret General; detailed rules remain in the following subclauses.
\item[13.7.6.2] \textbf{Matching of partial specializations.} Defines the template context for Matching of partial specializations; track dependence, lookup time, substitution, constraints, instantiation, and specialization ordering separately.
\item[13.7.6.3] \textbf{Partial ordering of partial specializations.} Defines the template context for Partial ordering of partial specializations; track dependence, lookup time, substitution, constraints, instantiation, and specialization ordering separately.
\item[13.7.6.4] \textbf{Members of class template partial specializations.} Defines the template context for Members of class template partial specializations; track dependence, lookup time, substitution, constraints, instantiation, and specialization ordering separately.
\item[13.7.7.1] \textbf{General.} Frames the terminology, applicability, and relationships needed to interpret General; detailed rules remain in the following subclauses.
\item[13.7.7.2] \textbf{Function template overloading.} Contributes to overload selection for Function template overloading: construct the candidate set, test viability, rank conversion sequences, apply tie-breakers, and require a unique best result.
\item[13.7.7.3] \textbf{Partial ordering of function templates.} Defines the template context for Partial ordering of function templates; track dependence, lookup time, substitution, constraints, instantiation, and specialization ordering separately.
\end{description}
\Needspace{8\baselineskip}
\subsection*{13.8 Name resolution}
\begin{longtable}{@{}>{\RaggedRight\arraybackslash\columncolor{CornellReferenceCueGray}}p{0.28\textwidth}>{\RaggedRight\arraybackslash}p{0.64\textwidth}@{}}
\rowcolor{CornellReferenceNavy}\color{white}\textbf{Cue / recall prompt} & \color{white}\textbf{Notes}\\\endfirsthead
\rowcolor{CornellReferenceNavy}\color{white}\textbf{Cue / recall prompt} & \color{white}\textbf{Notes (continued)}\\\endhead
\arrayrulecolor{CornellReferenceLineGray}
\textbf{What rule applies to General?}\\[-1mm]{\scriptsize\CornellClauseRef{13.8.1}} & Frames the terminology, applicability, and relationships needed to interpret General; detailed rules remain in the following subclauses.\\\hline
\textbf{What rule applies to Locally declared names?}\\[-1mm]{\scriptsize\CornellClauseRef{13.8.2}} & Defines the core-language rules for Locally declared names, including formation, semantic effect, interactions with type and lookup, and any ill-formed or implementation-dependent cases.\\\hline
\textbf{What rule applies to Dependent names?}\\[-1mm]{\scriptsize\CornellClauseRef{13.8.3}} & Defines the template context for Dependent names; track dependence, lookup time, substitution, constraints, instantiation, and specialization ordering separately.\\\hline
\textbf{What rule applies to Dependent name resolution?}\\[-1mm]{\scriptsize\CornellClauseRef{13.8.4}} & Defines the template context for Dependent name resolution; track dependence, lookup time, substitution, constraints, instantiation, and specialization ordering separately.\\\hline
\end{longtable}
\begin{CornellReferenceSummaryBox}Name resolution supplies the core-language semantics portion of this clause. The key review move is to connect its local wording to the clause-wide model, then classify each condition by normative force before relying on it.\end{CornellReferenceSummaryBox}
\Needspace{5\baselineskip}\paragraph{Detailed coverage ledger.}
\begin{description}[style=nextline,leftmargin=2.8cm,font=\normalfont\bfseries\color{CornellReferenceTeal}]
\item[13.8.3.1] \textbf{General.} Frames the terminology, applicability, and relationships needed to interpret General; detailed rules remain in the following subclauses.
\item[13.8.3.2] \textbf{Dependent types.} Defines the template context for Dependent types; track dependence, lookup time, substitution, constraints, instantiation, and specialization ordering separately.
\item[13.8.3.3] \textbf{Type-dependent expressions.} Defines the template context for Type-dependent expressions; track dependence, lookup time, substitution, constraints, instantiation, and specialization ordering separately.
\item[13.8.3.4] \textbf{Value-dependent expressions.} Defines the template context for Value-dependent expressions; track dependence, lookup time, substitution, constraints, instantiation, and specialization ordering separately.
\item[13.8.3.5] \textbf{Dependent template arguments.} Defines the template context for Dependent template arguments; track dependence, lookup time, substitution, constraints, instantiation, and specialization ordering separately.
\item[13.8.4.1] \textbf{Point of instantiation.} Defines the core-language rules for Point of instantiation, including formation, semantic effect, interactions with type and lookup, and any ill-formed or implementation-dependent cases.
\item[13.8.4.2] \textbf{Candidate functions.} Contributes to overload selection for Candidate functions: construct the candidate set, test viability, rank conversion sequences, apply tie-breakers, and require a unique best result.
\end{description}
\Needspace{8\baselineskip}
\subsection*{13.9 Template instantiation and specialization}
\begin{longtable}{@{}>{\RaggedRight\arraybackslash\columncolor{CornellReferenceCueGray}}p{0.28\textwidth}>{\RaggedRight\arraybackslash}p{0.64\textwidth}@{}}
\rowcolor{CornellReferenceNavy}\color{white}\textbf{Cue / recall prompt} & \color{white}\textbf{Notes}\\\endfirsthead
\rowcolor{CornellReferenceNavy}\color{white}\textbf{Cue / recall prompt} & \color{white}\textbf{Notes (continued)}\\\endhead
\arrayrulecolor{CornellReferenceLineGray}
\textbf{What rule applies to General?}\\[-1mm]{\scriptsize\CornellClauseRef{13.9.1}} & Frames the terminology, applicability, and relationships needed to interpret General; detailed rules remain in the following subclauses.\\\hline
\textbf{What rule applies to Implicit instantiation?}\\[-1mm]{\scriptsize\CornellClauseRef{13.9.2}} & Defines the core-language rules for Implicit instantiation, including formation, semantic effect, interactions with type and lookup, and any ill-formed or implementation-dependent cases.\\\hline
\textbf{What rule applies to Explicit instantiation?}\\[-1mm]{\scriptsize\CornellClauseRef{13.9.3}} & Defines the core-language rules for Explicit instantiation, including formation, semantic effect, interactions with type and lookup, and any ill-formed or implementation-dependent cases.\\\hline
\textbf{What rule applies to Explicit specialization?}\\[-1mm]{\scriptsize\CornellClauseRef{13.9.4}} & Defines the template context for Explicit specialization; track dependence, lookup time, substitution, constraints, instantiation, and specialization ordering separately.\\\hline
\end{longtable}
\begin{CornellReferenceSummaryBox}Template instantiation and specialization supplies the core-language semantics portion of this clause. The key review move is to connect its local wording to the clause-wide model, then classify each condition by normative force before relying on it.\end{CornellReferenceSummaryBox}
\Needspace{8\baselineskip}
\subsection*{13.10 Function template specializations}
\begin{longtable}{@{}>{\RaggedRight\arraybackslash\columncolor{CornellReferenceCueGray}}p{0.28\textwidth}>{\RaggedRight\arraybackslash}p{0.64\textwidth}@{}}
\rowcolor{CornellReferenceNavy}\color{white}\textbf{Cue / recall prompt} & \color{white}\textbf{Notes}\\\endfirsthead
\rowcolor{CornellReferenceNavy}\color{white}\textbf{Cue / recall prompt} & \color{white}\textbf{Notes (continued)}\\\endhead
\arrayrulecolor{CornellReferenceLineGray}
\textbf{What rule applies to General?}\\[-1mm]{\scriptsize\CornellClauseRef{13.10.1}} & Frames the terminology, applicability, and relationships needed to interpret General; detailed rules remain in the following subclauses.\\\hline
\textbf{What rule applies to Explicit template argument specification?}\\[-1mm]{\scriptsize\CornellClauseRef{13.10.2}} & Defines the template context for Explicit template argument specification; track dependence, lookup time, substitution, constraints, instantiation, and specialization ordering separately.\\\hline
\textbf{What rule applies to Template argument deduction?}\\[-1mm]{\scriptsize\CornellClauseRef{13.10.3}} & Defines the template context for Template argument deduction; track dependence, lookup time, substitution, constraints, instantiation, and specialization ordering separately.\\\hline
\textbf{What rule applies to Overload resolution?}\\[-1mm]{\scriptsize\CornellClauseRef{13.10.4}} & Contributes to overload selection for Overload resolution: construct the candidate set, test viability, rank conversion sequences, apply tie-breakers, and require a unique best result.\\\hline
\end{longtable}
\begin{CornellReferenceSummaryBox}Function template specializations supplies the operation semantics portion of this clause. The key review move is to connect its local wording to the clause-wide model, then classify each condition by normative force before relying on it.\end{CornellReferenceSummaryBox}
\Needspace{5\baselineskip}\paragraph{Detailed coverage ledger.}
\begin{description}[style=nextline,leftmargin=2.8cm,font=\normalfont\bfseries\color{CornellReferenceTeal}]
\item[13.10.3.1] \textbf{General.} Frames the terminology, applicability, and relationships needed to interpret General; detailed rules remain in the following subclauses.
\item[13.10.3.2] \textbf{Deducing template arguments from a function call.} Defines the template context for Deducing template arguments from a function call; track dependence, lookup time, substitution, constraints, instantiation, and specialization ordering separately.
\item[13.10.3.3] \textbf{Deducing template arguments taking the address of a function template.} Defines the template context for Deducing template arguments taking the address of a function template; track dependence, lookup time, substitution, constraints, instantiation, and specialization ordering separately.
\item[13.10.3.4] \textbf{Deducing conversion function template arguments.} Specifies source and destination conditions for Deducing conversion function template arguments, the resulting type or value category, and any information-loss, pointer, qualification, or lifetime consequence.
\item[13.10.3.5] \textbf{Deducing template arguments during partial ordering.} Defines the template context for Deducing template arguments during partial ordering; track dependence, lookup time, substitution, constraints, instantiation, and specialization ordering separately.
\item[13.10.3.6] \textbf{Deducing template arguments from a type.} Defines the template context for Deducing template arguments from a type; track dependence, lookup time, substitution, constraints, instantiation, and specialization ordering separately.
\item[13.10.3.7] \textbf{Deducing template arguments from a function declaration.} Defines the template context for Deducing template arguments from a function declaration; track dependence, lookup time, substitution, constraints, instantiation, and specialization ordering separately.
\end{description}
\section{Language-Lawyer and Library-Usage Pitfalls}
\begin{CornellReferencePitfallBox}\begin{itemize}[label=\textcolor{CornellReferenceAmber}{$\blacktriangleright$}]
\item Performing non-dependent lookup at instantiation time.
\item Treating every substitution error as non-diagnostic removal.
\item Assuming logically equivalent constraints necessarily have identical normalized structure.
\item Specializing or instantiating in an impermissible context.
\item Confusing function-template overloading with class-template partial specialization.
\item Treating a note, example, or recommended practice as though it imposed a mandatory program rule.
\end{itemize}\end{CornellReferencePitfallBox}
\section{Programmer and Implementation Checklist}
\begin{CornellReferenceChecklistBox}\begin{itemize}[label=$\square$]
\item Classify every parameter and validate its argument.
\item Identify dependent names, types, expressions, and template arguments.
\item Track definition and instantiation lookup contexts.
\item Normalize constraints before reasoning about subsumption.
\item Separate deduction, substitution, viability, and overload ranking.
\item Check implicit and explicit instantiation timing.
\item Apply specialization matching and partial ordering precisely.
\item For \CornellClauseRef{13.1}, verify preamble against the precise rule category and all cited dependencies.
\item For \CornellClauseRef{13.2}, verify template parameters against the precise rule category and all cited dependencies.
\item For \CornellClauseRef{13.3}, verify names of template specializations against the precise rule category and all cited dependencies.
\item For \CornellClauseRef{13.4}, verify template arguments against the precise rule category and all cited dependencies.
\end{itemize}\end{CornellReferenceChecklistBox}
\section{Clause Synthesis}
\begin{CornellReferenceOrientationBox}Defines template parameters and arguments, constraints, declaration forms, specialization, dependent names, instantiation, deduction, partial ordering, and overload resolution for templates. Provides compile-time abstraction and controlled generation of declarations and definitions. A sound review distinguishes syntax and participation from runtime preconditions, keeps notes and examples non-mandatory, and follows cross-clause dependencies before making a portability claim.\end{CornellReferenceOrientationBox}
\section{Self-Test Questions}
\begin{enumerate}
\item What role does Preamble play in the clause's overall model?
\item Which normative distinctions must be kept separate when applying Template parameters?
\item How would you audit a program or implementation that depends on Names of template specializations?
\item Compare Template arguments with Template constraints; what different question does each answer?
\item What role does Template constraints play in the clause's overall model?
\item Which normative distinctions must be kept separate when applying Type equivalence?
\item How would you audit a program or implementation that depends on Template declarations?
\item Compare Name resolution with Template instantiation and specialization; what different question does each answer?
\item What role does Template instantiation and specialization play in the clause's overall model?
\item Which normative distinctions must be kept separate when applying Function template specializations?
\item How would you audit a program or implementation that depends on General?
\item Compare Template type arguments with Template non-type arguments; what different question does each answer?
\item What role does Template non-type arguments play in the clause's overall model?
\item Which normative distinctions must be kept separate when applying Template template arguments?
\item Scenario: a program relies on an unstated implementation behavior while using General. What must be checked before calling the program portable?
\item Scenario: a program relies on an unstated implementation behavior while using Constraints. What must be checked before calling the program portable?
\end{enumerate}
\clearpage\section{Answer Key}
\begin{enumerate}
\item Preamble is one part of the clause's core-language semantics model. Its role is understood by connecting the local rule in 13.1 to the clause orientation and its dependent concepts.
\item Keep formation and well-formedness separate from runtime preconditions and effects. Also distinguish mandatory wording from notes, implementation choice from undefined behavior, and interface declarations from detailed contracts; see 13.2.
\item Audit Names of template specializations by locating the controlling wording in 13.3, listing inputs and type constraints, proving any preconditions, recording effects and failure behavior, and checking lifetime, invalidation, complexity, and synchronization where applicable.
\item Template arguments and Template constraints occupy different steps or facility groups. Analyze each under its own subclause, then follow the explicit dependency between them rather than transferring one set of guarantees to the other; begin at 13.4.
\item Template constraints is one part of the clause's requirements model. Its role is understood by connecting the local rule in 13.5 to the clause orientation and its dependent concepts.
\item Keep formation and well-formedness separate from runtime preconditions and effects. Also distinguish mandatory wording from notes, implementation choice from undefined behavior, and interface declarations from detailed contracts; see 13.6.
\item Audit Template declarations by locating the controlling wording in 13.7, listing inputs and type constraints, proving any preconditions, recording effects and failure behavior, and checking lifetime, invalidation, complexity, and synchronization where applicable.
\item Name resolution and Template instantiation and specialization occupy different steps or facility groups. Analyze each under its own subclause, then follow the explicit dependency between them rather than transferring one set of guarantees to the other; begin at 13.8.
\item Template instantiation and specialization is one part of the clause's core-language semantics model. Its role is understood by connecting the local rule in 13.9 to the clause orientation and its dependent concepts.
\item Keep formation and well-formedness separate from runtime preconditions and effects. Also distinguish mandatory wording from notes, implementation choice from undefined behavior, and interface declarations from detailed contracts; see 13.10.
\item Audit General by locating the controlling wording in 13.10.3.1, listing inputs and type constraints, proving any preconditions, recording effects and failure behavior, and checking lifetime, invalidation, complexity, and synchronization where applicable.
\item Template type arguments and Template non-type arguments occupy different steps or facility groups. Analyze each under its own subclause, then follow the explicit dependency between them rather than transferring one set of guarantees to the other; begin at 13.4.2.
\item Template non-type arguments is one part of the clause's core-language semantics model. Its role is understood by connecting the local rule in 13.4.3 to the clause orientation and its dependent concepts.
\item Keep formation and well-formedness separate from runtime preconditions and effects. Also distinguish mandatory wording from notes, implementation choice from undefined behavior, and interface declarations from detailed contracts; see 13.4.4.
\item First identify the exact requirement in 13.10.3.1, then classify the relied-on behavior as required, unspecified, implementation-defined, conditionally supported, or undefined. Portability requires satisfying the program obligations and avoiding dependence on a choice not guaranteed in the target environments.
\item First identify the exact requirement in 13.5.2, then classify the relied-on behavior as required, unspecified, implementation-defined, conditionally supported, or undefined. Portability requires satisfying the program obligations and avoiding dependence on a choice not guaranteed in the target environments.
\end{enumerate}
\section{Key-Term Recap}
\begin{longtable}{@{}>{\RaggedRight\arraybackslash}p{0.28\textwidth}>{\RaggedRight\arraybackslash}p{0.64\textwidth}@{}}\toprule\textbf{Term} & \textbf{Working meaning}\\\midrule\endfirsthead\toprule\textbf{Term} & \textbf{Working meaning (continued)}\\\midrule\endhead\bottomrule\endfoot
dependent name & A name whose interpretation depends on template arguments.\\
instantiation & Formation of a specialization from a template and arguments.\\
specialization & A template entity associated with a particular argument pattern.\\
constraint normalization & Transformation of constraint expressions into the form used for identity and subsumption.\\
substitution failure & A failure in the immediate substitution context that can remove a candidate rather than diagnose the program.\\
point of instantiation & A context governing lookup and semantic formation of an instantiated specialization.\\
partial ordering & Rules that determine which template or specialization is more specialized.\\
\end{longtable}
\CornellTakeaway{Template semantics are context-sensitive: dependence, lookup time, substitution, constraints, instantiation, and ordering must be analyzed in sequence.}
\end{document}
